\section{Degenerate slow spaces and a sharp bound on strong preparations}\label{sec:degenerateextension}
A scalar overlap does not classify a degenerate slow eigenspace. There is,
however, a complete finite test for a specified asymptotic distance, and a
simple sharp bound on strong preparations that works for any number of states.

Let $S=D^{-1/2}MD^{1/2}$, and let $P$ be the orthogonal projector onto the
entire slowest nonstationary eigenspace of dimension $r$. Put
\[
H=D^{-1/2}PD^{-1/2},\qquad
Q(\beta)=\pi(\beta)^\top H\pi(\beta).
\]
Since $P\sqrt{\pi_b}=0$, $H\pi_b=0$. Spectral expansion gives
\[
\sum_i\frac{(p_i(t;\beta)-\pi_{bi})^2}{\pi_{bi}}
=e^{2\lambda_2t}Q(\beta)+o(e^{2\lambda_2t}).
\]
Relative entropy has the same leading term multiplied by $1/2$. Thus strict
ordering of $Q$ certifies asymptotic weighted Euclidean and relative-entropy
reversals. If two specified preparations have equal $Q$, subleading terms can
decide their eventual ordering; the leading coefficient alone does not decide
that pair. The existence test below remains complete over all allowed initial
temperatures: if $Q'$ has no sign change on a side, analyticity and the bath
zero make $Q$ strictly ordered there. For other norms a degenerate projection can give different orderings;
no norm-independent scalar assertion is made. A strong preparation is precisely
$Q=0$ away from the bath, equivalently $PD^{-1/2}\pi(\beta)=0$.

\begin{theorem}[State count and slow-space dimension]\label{thm:strongdimension}
For $n$ distinct energies and a slow eigenspace of dimension $r$, there are at
most $n-r-1$ distinct nonbath strong Gibbs preparations. For every prescribed
set of $n-r-1$ distinct positive nonbath inverse temperatures, and every
positive finite bath, this bound is attained by a strictly positive reversible
generator with an $r$-dimensional slow eigenspace.
\end{theorem}
\begin{proof}
Any $k\leq n$ distinct Gibbs vectors are linearly independent. To see this,
take their first $k$ coordinates and multiply each column by its positive
partition function. The resulting evaluation matrix of $k$ distinct
exponentials has nonzero determinant by the exponential Chebyshev property.
Multiplication by $D^{-1/2}$ preserves independence.

The kernel of $P$ has dimension $n-r$ and already contains the transformed
bath vector $\sqrt{\pi_b}$. It therefore contains at most $n-r-1$ additional
Gibbs vectors. For attainability, let $U$ be the span of the transformed bath
and the $n-r-1$ prescribed transformed Gibbs vectors. They are independent,
so $\dim U=n-r$. Let $P$ project orthogonally onto $U^\perp$, and define
\[
S=-I+ss^\top+\epsilon P,\qquad s=\sqrt{\pi_b},\qquad
M=D^{1/2}SD^{-1/2}.
\]
Choose $0<\epsilon<1$ sufficiently small that
$s_is_j+\epsilon P_{ij}>0$ for every $i\ne j$; this is possible because all
$s_i$ are positive and there are only finitely many entries. Then $M$ has
positive off-diagonal entries, vanishing column sums, detailed balance, and
spectrum $0$, $-1+\epsilon$ with multiplicity $r$, and $-1$ with multiplicity
$n-r-1$. Its strong preparations include all the prescribed ones, and the
bound just proved excludes any others.
\end{proof}

For four states the maximum number of distinct strong preparations is therefore
two, one, or zero when the slow multiplicity is one, two, or three. This bound
counts exact absence of the \emph{whole} slow projection; canceling one chosen
vector in a degenerate eigenspace is insufficient. It also clarifies why a
classification must retain spectral multiplicity. The realization proof extends the rank-one perturbation in Klich et al.,
Appendix VIII, to an orthogonal projector of rank $r$. No novelty claim is made
for the underlying linear-independence or projector-construction principles;
publication priority for the sharp Gibbs-preparation bound remains to be checked.
For $r=1$, the $n-2$ nonbath-zero bound already follows explicitly from
Emmerson's Theorem 7.3 and Corollary 7.5, as detailed in
Section~\ref{sec:novelty}; no new scalar strong-temperature bound is claimed.
The dimension argument here also covers simultaneous cancellation when $r>1$.

\paragraph{Finite classification when $r=2$.}
Write $f_i=e^{-\beta E_i}$ and $Z=\sum_i f_i$. Then
\[
Q'=\frac{R(\beta)}{Z^3},\qquad
R=(f^\top Hf)'Z-2(f^\top Hf)Z'.
\]
For four states, $R$ contains at most sixteen distinct exponential terms.
There are twelve terms $f_i^2f_k$ with $i\ne k$, with coefficients
$2(E_k-E_i)(H_{ii}-H_{ik})$, and four terms with distinct indices $i,j,k$,
with coefficients
\[
2\{H_{ij}(2E_k-E_i-E_j)+H_{ik}(2E_j-E_i-E_k)
+H_{jk}(2E_i-E_j-E_k)\}.
\]
All four pure cubes have zero coefficient; coincident exponent sums can only
reduce the count. After combining equal exponents, a nonzero $m$-term $R$
has at most $m-1\leq15$ real zeros counting multiplicity. It cannot vanish
identically: that would make $Q$ constantly zero since $Q(\beta_b)=0$, forcing
$P$ to annihilate all Gibbs vectors, which span the full state space.

Isolate these finitely many zeros on $(0,\beta_b)$ and $(\beta_b,\infty)$,
then evaluate the sign of $R$ once on each intervening interval. An ordinary
asymptotic cooling effect exists exactly when some cooling interval has $R>0$;
a heating effect exists exactly when some heating interval has $R<0$.
This follows directly by differentiating the nonnegative amplitude $Q$.
Exponential root isolation can be organized by the repeated Rolle procedure:
factor out the smallest exponent, differentiate to reduce the number of terms
by one, recursively isolate derivative roots, and use monotonicity between
those roots. Limiting endpoint signs follow from the leading nonzero term.
Exact tangencies must be retained. This is a finite mathematical prescription;
it does not assert that floating-point sampling or arbitrary transcendental
equality testing is a certified terminating decision algorithm.

For $r=3$, $S=-\mu(I-ss^\top)$ for some $\mu>0$, hence
$M=\mu(\pi_b\mathbf1^\top-I)$. This reset generator is stochastically monotone,
so the earlier all-time convex-divergence ordering applies. There is no
asymptotic cooling or heating effect in the distances just specified.
Together with Theorem~\ref{thm:finiteclassification}, these cases cover all
slow multiplicities of irreducible reversible four-state systems for
asymptotic relative entropy and weighted Euclidean distance. The arguments
also allow zero transition rates when irreducibility and the stated spectral
conditions remain valid; disconnected generators require a different target
equilibrium and are outside this thermalizing model.
